\chapter{Șiruri și serii de funcții}

\section{Convergență punctuală și convergență uniformă}

\begin{definition}\label{def:conv-punct}
  \index{șiruri funcții!convergență!punctuală}
  Fie $(X, d) $ un spațiu metric și $ f_n : X \to \RR $ termenul general
  al unui șir de funcții.

  Fie $ f : X \to \RR $ o funcție arbitrară.

  Spunem că șirul $ (f_n) $ \emph{converge punctual} (\emph{simplu}) la $ f $
  dacă are loc:
  \[
    \lim_{n \to \infty} f_n(x) = f(x), \quad \forall x \in X.
  \]

  Notăm acest lucru cu $ f_n \xrar{PC} f $ și numim $ f $ \emph{limita punctuală}
  a șirului $ (f_n) $.
\end{definition}

Celălalt tip de convergență care ne va interesa este definit mai jos:
\begin{definition}\label{def:conv-unif}
  \index{șiruri funcții!convergență!uniformă}
  În condițiile și cu notațiile de mai sus, spunem că șirul $ (f_n) $ este
  \emph{uniform convergent} la $ f $ dacă:
  \[
    \forall \varepsilon > 0, \exists N_\varepsilon > 0 \text{ a.î. } %
    |f_n(x) - f(x)| < \varepsilon, \forall n \geq N_\varepsilon, \forall x \in X.
  \]

  Vom nota această situație cu $ f_n \xrar{UC} f $.
\end{definition}

În exerciții se va folosi mai mult caracterizarea:
\begin{proposition}\label{pr:conv-unif}
  În condițiile și cu notațiile de mai sus, șirul $(f_n)$ converge
  uniform la $ f $ dacă și numai dacă:
  \[
    \lim_{n \to \infty} \sup_{x \in X} |f_n(x) - f(x)| = 0.
  \]
\end{proposition}

Legătura între cele două tipuri de convergență este dată de:
\begin{theorem}\label{thm:UC-PC}
  Orice șir de funcții uniform convergent pe un interval este punctual
  convergent pe același interval.
\end{theorem}

Reciproca este falsă, după cum arată contraexemplul: fie $ [a, b] = [0,1] $
și definim șirul de funcții $ f_n(x) = x^n, n \geq 1 $.

Pentru orice $ x \in [a, b] $, avem:
\[
  \lim_{n \to \infty} f_n(x) = %
  \begin{cases}
    0, & x \in [0, 1) \\
    1, & x = 1
  \end{cases}
\]

Rezultă că $ f_n \xrar{PC} f $, unde $ f $ este funcția definită pe
cazuri de limita de mai sus. Dar se poate vedea imediat că
\[
  \lim_{n \to \infty} \sup_{x \in X} |f_n(x) - f(x)| = 1 \neq 0,
\]
deci șirul este doar punctual convergent, nu uniform convergent.

%%%%%%%%%%%%%%%%%%%%%%%%%%%%%%%%%%%%%%%%%%%%%%%%%%%%%%%%%%%%%%%%%%%%%%
\section{Transferul proprietăților}

Unele dintre proprietățile funcțiilor sînt transferate de la termenii șirurilor
la funcția-limită, iar altele nu, aceasta oferindu-ne uneori metode de calcul,
iar alteori, metode de demonstrație.

\begin{theorem}[Transfer de continuitate]\label{thm:transf-cont}
  \index{șiruri funcții!transfer de continuitate}
  Dacă $ f_n $ sînt funcții continue, iar șirul $ f_n $ converge uniform
  la $ f $, atunci funcția $ f $ este continuă.
\end{theorem}

Rezultă de aici că avem o metodă de a demonstra că nu are loc
continuitatea uniformă: dacă $ f_n $ sînt funcții continue, iar $ f $,
obținută din convergența punctuală, \emph{nu} este continuă, rezultă
că $ f_n $ \emph{nu} tinde uniform la $ f $.

\begin{theorem}[Integrare termen cu termen]\label{thm:int-termen}
  \index{șiruri funcții!integrare termen cu termen}
  Fie $ f_n, f : [a, b] \to \RR $ funcții continue. Dacă $ f_n $
  converge uniform la $ f $, atunci are loc proprietatea de
  integrare termen cu termen, adică:
  \[
    \lim_{n \to \infty} \int_a^b f_n(x) dx = \int_a^b f(x) dx.
  \]
\end{theorem}

\begin{theorem}[Derivare termen cu termen]\label{thm:derivare-termen}
  \index{șiruri funcții!derivare termen cu termen}
  Presupunem că funcțiile $ f_n $ sînt derivabile, pentru orice $ n \in \NN $.
  Dacă șirul $ f_n $ converge punctual la $ f $ și dacă există funcția
  $ g : [a, b] \to \RR $ astfel încît $ f'_n $ converge uniform la $ g $,
  atunci $ f $ este derivabilă și $ f' = g $.
\end{theorem}

%%%%%%%%%%%%%%%%%%%%%%%%%%%%%%%%%%%%%%%%%%%%%%%%%%%%%%%%%%%%%%%%%%%%%%

\section{Serii de funcții}

Pentru convergența seriilor de funcții, avem un singur criteriu de utilizat:
\begin{theorem}[Weierstrass]\label{thm:wei-serii-f}
  \index{serii!funcții!Weierstrass}
  Dacă există un șir cu termeni pozitivi $ a_n $ astfel încît $ |u_n(x)| \leq a_n $
  pentru orice $ x \in X $, iar seria $ \sum a_n $ converge, atunci seria
  $ \sum u_n $ converge uniform.
\end{theorem}

Pentru proprietățile de transfer, avem:
\begin{theorem}\label{thm:transfer}
  \begin{itemize}
  \item \emph{Transfer de continuitate:} Dacă $ u_n $ sînt funcții continue,
    iar seria $ \sum u_n $ converge uniform la $ f $, atunci funcția $ f $
    este continuă.
  \item \emph{Integrare termen cu termen}: Dacă seria $ \sum u_n $ converge
    uniform la $ f $, atunci $ f $ este integrabilă și avem:
    \[
      \int_a^b \sum_n u_n(x) dx = \sum_n \int_a^b u_n(x) dx.
    \]
  \item \emph{Derivare termen cu termen:} Presupunem că toate funcțiile $ u_n $
    sînt derivabile. Dacă seria $ \sum u_n $ converge punctual la $ f $ și
    dacă există $ g : [a,b] \to \RR $ astfel încît $ \sum_n u'_n $
    converge uniform la $ g $, atunci $ f $ este derivabilă și $ f' = g $.
  \end{itemize}
\end{theorem}

%%%%%%%%%%%%%%%%%%%%%%%%%%%%%%%%%%%%%%%%%%%%%%%%%%%%%%%%%%%%%%%%%%%%%%

\section{Polinomul Taylor și seria Taylor}

Orice funcție cu anumite proprietăți poate fi aproximată cu un polinom:

\begin{definition}\label{def:pol-taylor}
  \index{serii!polinom Taylor}
  Fie $ I \seq \RR $ un interval deschis și $ f: I \to \RR $ o funcție de clasă
  $ \kal{C}^m(I) $. Pentru orice $ a \in I $, definim \emph{polinomul Taylor}
  de gradul $ n \leq m $ asociat funcției $ f $ în punctual $ a $ prin:
  \[
    T_{n,f,a}(x) = \sum_{k = 0}^n \frac{f^{(k)}(a)}{k!} (x - a)^k.
  \]
  \emph{Restul} (eroarea de aproximare) este definit prin:
  \[
    R_{n,f,a} = f(x) - T_{n,f,a}(x).
  \]
\end{definition}

Acest polinom poate fi mai departe utilizat pentru a studia \emph{seria Taylor}
asociată unei funcții.

\begin{theorem}\label{thm:seria-taylor}
  \index{serii!Taylor}
  Fie $ a < b $ și $ f \in \kal{C}^\infty([a,b]) $ astfel încît să existe
  $ M > 0 $ cu proprietatea că $ \forall n \in \NN $ și $ x \in [a, b] $, avem
  $ |f^{(n)}(x)| \leq M $.

  Atunci pentru orice $ x_0 \in (a,b) $, seria Taylor a lui $ f $ în jurul
  punctului $ x_0 $ este uniform convergentă pe $ [a,b] $ și suma ei este
  funcția $ f $, adică avem:
  \[
    f(x) = \sum_{n \geq 0} \frac{f^{(n)}(x_0)}{n!} (x - x_0)^n, \quad %
    \forall x \in [a,b].
  \]
\end{theorem}

Pentru cazul particular $ x_0 = 0 $, seria se numește \emph{Maclaurin}.
\index{serii!Maclaurin}

%%%%%%%%%%%%%%%%%%%%%%%%%%%%%%%%%%%%%%%%%%%%%%%%%%%%%%%%%%%%%%%%%%%%%%

\section{Serii de puteri}

Seriile de puteri sînt un caz particular al seriilor de funcții, luînd doar
funcții de tip polinomial.

\begin{definition}\label{def:serii-puteri}
  \index{serii!puteri}
  Fie $ (a_n) $ un șir de numere reale și fie $ a \in \RR $.

  Seria $ \sum_{n \geq 0} a_n(x - a)^n $ se numește \emph{seria de puteri %
  centrată în $ a $}, definită de șirul $ a_n $.
\end{definition}

Toate rezultatele privitoare la serii de funcții sînt valabile
și pentru serii de puteri. Rezultatele specifice urmează.

\begin{theorem}[Abel]\label{thm:abel-s-put}
  \index{serii!puteri!Abel}
  Pentru orice serie de puteri $ \sum a_n x^n $ există un număr
  $ 0 \leq R \leq \infty $ astfel încît:
  \begin{itemize}
  \item Seria este absolut convergentă pe intervalul $ (-R, R) $;
  \item Seria este divergentă pentru orice $ |x| > R $;
  \item Seria este uniform convergentă pe $ [-r, r] $, unde $ 0 < r < R $.
  \end{itemize}

  Numărul $ R $ se numește \emph{raza de convergență} a seriei, iar
  intervalul $ (-R, R) $ se numește \emph{intervalul de convergență}.
\end{theorem}

Calculul razei de convergență se poate face cu unul dintre următoarele criterii:

\begin{theorem}[Cauchy-Hadamard]\label{thm:c-h-raza}
  \index{serii!puteri!Cauchy-Hadamard}
  Fie $ \sum a_n x^n $ o serie de puteri, $ R $ raza sa de
  convergență și definim:
  \[
    \omega = \lim \sup \sqrt[n]{|a_n|}.
  \]
  Atunci:
  \begin{itemize}
  \item $ R = \omega^{-1} $ dacă $ 0 < \omega < \infty $;
  \item $ R = 0 $ dacă $ \omega = \infty $;
  \item $ R = \infty $ dacă $ \omega = 0 $.
  \end{itemize}
\end{theorem}

\begin{theorem}\label{thm:raport-raza}
  \index{serii!puteri!raport}
  Raza de convergență se poate calcula cu formula:
  \[
    R = \lim_{n \to \infty} \frac{|a_n|}{|a_{n+1}|}.
  \]
\end{theorem}

\begin{remark}\label{rk:transfer-puteri}
  Din natura seriilor de puteri, teoremele de derivare și integrare termen
  cu termen sînt automate. Așadar, dacă $ \sum a_n (x - a)^n $ este o
  serie de puteri iar $ S(x) $ este suma sa, atunci:
  \begin{itemize}
  \item Seria derivatelor, $ \sum na_n(x - a)^{n-1} $, are aceeași rază de
    convergență cu seria inițială, iar suma sa este $ S'(x) $;
  \item Seria primitivelor, $ \sum a_n \dfrac{(x - a)^{n+1}}{n+1} $,
    are aceeași rază de convergență cu seria inițială, iar suma sa este o
    primitivă a lui $ S $.
  \end{itemize}
\end{remark}

%%%%%%%%%%%%%%%%%%%%%%%%%%%%%%%%%%%%%%%%%%%%%%%%%%%%%%%%%%%%%%%%%%%%%%

\section{Exerciții}

1. Să se studieze convergența punctuală și uniformă a șirurilor de funcții:
\begin{enumerate}[(a)]
\item $ f_n : (0, 1) \to \RR, f_n(x) = \dfrac{1}{nx + 1}, n \geq 0 $;
\item $ f_n : [0, 1] \to \RR, f_n(x) = x^n - x^{2n}, n \geq 0 $;
\item $ f_n : \RR \to \RR, f_n(x) = \sqrt{x^2 + \dfrac{1}{n^2}}, n > 0 $;
\item $ f_n : [-1,1] \to \RR, f_n(x) = \dfrac{x}{1 + nx^2} $;
\item $ f_n : (-1, 1) \to \RR, f_n(x) = \dfrac{1 - x^n}{1 - x} $;
\item $ f_n : [0, 1] \to \RR, f_n(x) = \dfrac{2nx}{1 + n^2 x^2} $;
\item $ f_n : \RR_+ \to \RR, f_n(x) = \dfrac{x + n}{x + n + 1} $;
\item $ f_n : \RR_+ \to \RR, f_n(x) = \dfrac{x}{1 + nx^2} $.
\end{enumerate}

\vspace{1cm}

2. Să se arate că șirul de funcții dat de:
\[
  f_n : \RR \to \RR, \quad f_n(x) = \frac{1}{n} \arctan x^n
\]
converge uniform pe $ \RR $, dar:
\[
  \Big( \lim_{n \to \infty} f_n(x) \Big)'_{x = 1} \neq \lim_{n \to \infty} f'_n(1).
\]

Rezultatele diferă deoarece șirul derivatelor nu converge uniform pe $ \RR $.

\vspace{1cm}

3. Să se arate că șirul de funcții dat de:
\[
  f_n : [0, 1] \to \RR, \quad f_n(x) = nxe^{-nx^2}
\]
este convergent, dar:
\[
  \lim_{n \to \infty} \int_0^1 f_n(x) dx \neq \int_0^1 \lim_{n \to \infty} f_n(x) dx.
\]
Rezultatul se explică prin faptul că șirul nu este uniform convergent.
De exemplu, pentru $ x_n = \dfrac{1}{n} $, avem $ f_n(x_n) \to 1 $, dar,
în general, $ f_n(x) \to 0 $.

\vspace{1cm}

4. Să se dezvolte următoarele funcții în serie Maclaurin, precizînd și
domeniul de convergență:
\begin{enumerate}[(a)]
\item $ f(x) = e^x $;
\item $ f(x) = \sin x $;
\item $ f(x) = \cos x $;
\item $ f(x) = (1 + x)^\alpha, \alpha \in \RR $;
\item $ f(x) = \dfrac{1}{1 + x} $;
\item $ f(x) = \ln(1 + x) $;
\item $ f(x) = \arctan x $;
\item $ f(x) = \ln(1 + 5x)$;
\item $ f(x) = 3 \ln(2 + 3x) $.
\end{enumerate}

%%%%%%%%%%%%%%%%%%%%%%%%%%%%%%%%%%%%%%%%%%%%%%%%%%%%%%%%%%%%%%%%%%%%%%

5. Să se calculeze raza de convergență și mulțimea de convergență pentru
următoarele serii de puteri:
\begin{enumerate}[(a)]
\item $ \sum_{n \geq 0} x^n $;
\item $ \sum_{n \geq 1} n^n x^n $;
\item $ \sum_{n \geq 1} (-1)^{n+1} \dfrac{x^n}{n} $;
\item $ \sum_{n \geq 1} \dfrac{n^nx^n}{n!} $;
\item $ \sum \dfrac{(x - 1)^{2n}}{n \cdot 9^n} $;
\item $ \sum \dfrac{(x + 3)^n}{n^2} $.
\end{enumerate}

\vspace{1cm}

6. Găsiți mulțimea de convergență și suma seriei:
\[
  \sum_{n \geq 0} (-1)^n \frac{x^{2n+1}}{2n + 1}.
\]
\emph{Indicație}: Se derivează termen cu termen și rezultă seria
geometrică de rază $ -x^2 $, căreia i se poate calcula suma, care apoi
se integrează.

\vspace{1cm}

7. Să se calculeze cu o eroare mai mică decît $ 10^{-3}$ integralele:
\begin{enumerate}[(a)]
\item $ \ds\int_0^{\frac{1}{2}} \dfrac{\sin x}{x} dx $;
\item $ \ds\int_0^{\frac{1}{2}} \dfrac{\ln(1 + x)}{x} dx $;
\item $ \ds\int_0^1 e^{-x^2} dx $.
\end{enumerate}


\vspace{1cm}

8. Să se calculeze polinomul Taylor de grad 3 în jurul originii
pentru funcțiile:
\begin{enumerate}[(a)]
\item $ f(x) = 3 \ln(2 + x) $;
\item $ f(x) = \arctan x $;
\item $ f(x) = \sqrt{1 + 2x} $.
\end{enumerate}

%%% Local Variables:
%%% mode: latex
%%% TeX-master: "../am1-18-19"
%%% End:
